\documentclass[12pt,a4paper]{article}
\usepackage[T2A]{fontenc}
\usepackage[utf8]{inputenc}
\usepackage[russian]{babel}
\usepackage{amsmath,amssymb}

\begin{document}

\section{Решение квадратного уравнения}

\noindent\textit{Задача:} решить уравнение $2x^2+5x-12=0$.

\textit{Решение.} Рассмотрим квадратное уравнение общего вида:
\[
    ax^2+bx+c=0.
\]

В данной задаче $a=2$, $b=5$, $c=-12$.

Сначала вычислим дискриминант:
\[
    D=b^2-4ac=(5)^2-4\cdot 2\cdot(-12)=121.
\]

Поскольку дискриминант положителен ($D>0$), уравнение имеет два корня:
\[
    x_{1,2}=\frac{-b\pm\sqrt{D}}{2a}=\frac{-5\pm11}{4}.
\]

Следовательно, $x_1=\frac{-16}{4}=-4$, $x_2=\frac{6}{4}=\frac{3}{2}$.

\textit{Ответ:} $x_1=-4$, $x_2=\frac{3}{2}$.

\end{document}
